Una red de transporte sanitario de la Región Metropolitana debe decidir cada semana cuántos vehículos contratar, en qué bases ubicarlos y qué hubs habilitar, y cada día cómo atender entre 455 y 731 solicitudes semanales de muestras clínicas, fármacos refrigerados, insumos y equipos. Hoy esas decisiones se toman por separado y con reglas simples. Este informe propone una herramienta que las coordina.

\parrafo{Método.} La herramienta es una matheurística. Heurísticas GRASP y VND generan rutas diarias factibles; un modelo de particionamiento de conjuntos resuelto con Gurobi elige entre ellas la flota, las bases, los hubs y la forma de atender cada solicitud; y una generación de columnas con los precios duales del modelo agrega rutas que pueden bajar el costo. Se compara contra un caso base que representa la planificación por reglas: dimensiona la flota según la demanda e inserta cada solicitud en la ruta más barata. Un validador independiente revisa que cada solución cumpla todas las restricciones.

\parrafo{Resultados.} La matheurística es más barata que el caso base en las 80 semanas disponibles: un 9,9\,\% en las 40 semanas BIG (\$3,9 millones por semana) y un 15,7\,\% en las 40 SMALL. La diferencia es estadísticamente significativa y casi no depende de la semilla. El ahorro viene de rutear mejor: con casi la misma flota recorre un 16\,\% menos de kilómetros y externaliza un 28\,\% menos. Una semana BIG se resuelve en unos 25 minutos en un computador personal.

\parrafo{Diagnóstico.} El recurso escaso es el furgón frío: en 26 de las 40 semanas BIG se usan los 18 disponibles, y las solicitudes difíciles son insumos y equipos grandes con ventanas de una hora. La división de solicitudes entre vehículos es esencial (sin ella el costo sube un tercio), mientras que los hubs casi no aportan con estos datos.

\parrafo{Sensibilidad e infraestructura.} En 60 corridas de sensibilidad la ventaja se mantuvo en todos los escenarios. Lo que más mueve el costo es la demanda, que lo hace crecer más que proporcionalmente, y la congestión, que lo sube hasta un 11\,\%. Exigir puntualidad total no encarece la operación. Ninguna base es imprescindible, una décima base no paga su costo y la capacidad de las bases no limita; en cambio, dos furgones adicionales en las semanas de alta demanda bajan el costo y las externalizaciones.

\parrafo{Recomendaciones.} Se recomienda (1) implementar la herramienta en tres etapas, con una marcha blanca de cuatro semanas, un piloto de ocho semanas en las bases del poniente y el despliegue completo; (2) contratar dos furgones adicionales en forma flexible para las semanas de alta demanda; (3) negociar ventanas de atención más amplias con los centros que más externalizan, porque media hora más baja el costo cerca de un 4\,\%; (4) exigir puntualidad total, que no tiene costo; y (5) evaluar con más semanas el cierre de Talagante. Las principales limitaciones son que la solución es aproximada, que los tiempos de viaje suponen flujo libre y que la demanda se supone conocida al inicio de la semana.
